\section{conclusion}
Experimental results showed that the feasibility-only approach did not outperform ALM in either the Wall or Push-T environments. In the Wall environment, the feasibility-guided planner was able to generate actions. In Push-T, the feasibility model often optimized actions to be small or random, indicating limitations in the current feasibility-guided planning formulation. These results suggest that the current feasibility optimization approach may be insufficient for complex manipulation tasks that require effective exploration and robust constraint balancing.

Future work could investigate more robust constrained optimization methods, including Langevin-based Augmented Lagrangian approaches, to improve exploration of the optimization landscape while maintaining constraint satisfaction. In addition, feasibility learning may be more effective when used as an auxiliary regularizer rather than as a standalone objective, as suggested by strong performance in the Wall environment when combined with DINO-WM dynamics. Adaptive feasibility weighting schemes could further improve performance by dynamically balancing task and constraint terms during optimization. Finally, incorporating temperature-controlled noise during optimization may help escape poor local minima and reduce degenerate solutions, particularly in more complex environments such as Push-T.